% !TeX root = ../main.tex
\section{Reduction of Gajo's logarithmic constitutive relations}
\label{sec:gajo-elastic-equivalence}

A single reduction covers Gajo's elastic and poroplastic volumetric
responses.  We use the separated volumetric and isochoric energies and
logarithmic mineral law, with reversible mineral volume change as in
Appendix~C.1 of \citet{gajo2011}.  Skeleton plastic volume may evolve:
Gajo's \(J_p\) is our \(a^p\), so \(J^e=J/a^p\); the elastic case has
\(\mbf F^p=\mbf I\).  His mineral-volume ratio \(J_s\), pore pressure
\(p_w\), and reference solid fraction \(1-n_0\) correspond to
\(\bar J\), \(p\), and \(\phi_{s0}\), while \(K\) and \(K_s\) retain
their meanings.  We retain his elastic mineral factors \(J^e_{s-m}\) and
\(J^e_{s-f}\) for contact loading and pore pressure.  The comparison also
adopts his assumption that isotropic effective loading and pore pressure
induce equivalent isotropic mineral strains \citep[Sec.~3(c)]{gajo2010}.

In the present notation, Gajo's equations~(35)--(39), specialized with the
logarithmic skeleton energy in his equation~(47), give \citep{gajo2011}
\begin{align*}
 \frac13\tr\bs\tau^{\prime}
 &=K\left[\ln\left(\frac{J}{a^p}\right)-\ln J^e_{s-f}\right],
 % equation-id: eq:gajo-plastic-skeleton-stress
\\
 \bar J&=J^e_{s-m}J^e_{s-f},
 % equation-id: eq:gajo-plastic-mineral-factorization
\\
 K_s\ln J^e_{s-m}&=\frac{1}{3\phi_{s0}}\tr\bs\tau^{\prime},
 % equation-id: eq:gajo-plastic-contact-compression
\\
 K_s\ln J^e_{s-f}&=-p\bar J.
 % equation-id: eq:gajo-plastic-pressure-compression
\end{align*}
Eliminating the pressure factor from the skeleton stress gives
\begin{equation*}
 \frac13\tr\bs\tau^{\prime}
 =K\left[\ln\left(\frac{J}{a^p}\right)+\frac{p\bar J}{K_s}\right].
 % equation-id: eq:gajo-plastic-eliminated-stress
\end{equation*}
Taking the logarithm of the mineral factorization and substituting both
loading relations then yields
\begin{align*}
 K_s\ln\bar J
 &=\frac{1}{3\phi_{s0}}\tr\bs\tau^{\prime}-p\bar J,
 % equation-id: eq:gajo-plastic-volume-stress-elimination
\\
 &=\frac{K}{\phi_{s0}}\ln\left(\frac{J}{a^p}\right)
 +\frac{K}{\phi_{s0}K_s}p\bar J-p\bar J.
 % equation-id: eq:gajo-plastic-volume-pressure-elimination
\end{align*}
Consequently, the mineral state satisfies
\begin{equation}
 K_s\ln\bar J
 +\left(1-\frac{K}{\phi_{s0}K_s}\right)p\bar J
 -\frac{K}{\phi_{s0}}\ln\left(\frac{J}{a^p}\right)=0.
 \label{eq:gajo-plastic-scalar-mineral-closure}
\end{equation}
This is \eqref{eq:verification-mineral-factors}, including its dependence
on plastic distention.  Conversely, a positive root reconstructs the factors as
\(J^e_{s-f}=\exp(-p\bar J/K_s)\) and
\(J^e_{s-m}=\bar J/J^e_{s-f}\).  Substitution verifies all four factor
relations above, establishing equivalence in both directions on the
selected branch.  Setting \(\mbf F^p=\mbf I\) recovers the elastic
relations in Gajo's equations~(3.27), (3.32), and (3.34) \citep{gajo2010}.

The conversion from Kirchhoff to Cauchy stress uses total volume \(J\),
even though the logarithmic skeleton response depends on \(J/a^p\).
Subtracting pore pressure from the single-prime Cauchy stress gives
\begin{align*}
 \frac13\tr\bs\sigma
 &=\frac{K}{J}\ln\left(\frac{J}{a^p}\right)
 -\left(1-\frac{K\bar J}{JK_s}\right)p,
 % equation-id: eq:gajo-plastic-total-mean-stress
\\
 &=\frac13\tr\bs\sigma^{\prime\prime}_d-B_d p.
 % equation-id: eq:gajo-plastic-drained-mean-stress
\end{align*}
This agrees with Gajo's equations~(48) and (51) and the volumetric part of
\eqref{eq:total-stress-drained-pressure-split}.  Here the drained stress
retains the current plastic deformation, and \(B_d\) is evaluated using
the mineral volume at the current pressure.

To identify reversible pressure coupling, hold both total deformation and
plastic deformation fixed.  Differentiating
\eqref{eq:gajo-plastic-scalar-mineral-closure} gives
\begin{equation*}
 \left[\frac{K_s}{\bar J}
 +\left(1-\frac{K}{\phi_{s0}K_s}\right)p\right]
 \left.\pder{\bar J}{p}\right|_{\mbf F,\mbf F^p}
 =-\left(1-\frac{K}{\phi_{s0}K_s}\right)\bar J.
 % equation-id: eq:gajo-plastic-mineral-pressure-tangent
\end{equation*}
Substituting into the derivative of the total mean stress yields
\begin{align*}
 -\left.\pder{}{p}\left(\frac13\tr\bs\sigma\right)\right|_{\mbf F,\mbf F^p}
 &=B_d+p\left.\pder{B_d}{p}\right|_{\mbf F,\mbf F^p},
 % equation-id: eq:gajo-plastic-finite-tangent-relation
\\
 &=1-\frac{K\bar J}
 {J\left[K_s+\left(1-\dfrac{K}{\phi_{s0}K_s}\right)p\bar J\right]}.
 % equation-id: eq:gajo-plastic-biot-tangent-reduction
\end{align*}
The result is precisely \(B\) in \eqref{eq:poroplastic-biot-correction}.
Plastic history enters through the mineral state, while the derivative
measures a reversible perturbation with that history fixed.  The reduction therefore gives the same mineral equation, mean stresses,
and reversible pressure response under the stated assumptions.  Equivalence
of complete plastic loading paths would additionally require matching
yield surfaces, flow rules, and hardening laws.
